\section{Bốn mươi tám kỹ thuật}
\label{sec:ch10-vuon-bien-the}

\index{LIME!phân bố các biến thể}%
Đặt hai trục của hai mục trước cạnh nhau thì được bảng phân loại chính của khảo
sát, và bảng ấy có 48 dòng: 48 kỹ thuật phái sinh từ LIME, mỗi dòng mang một
hoặc nhiều chữ cái chỉ nhóm vấn đề, cùng các dấu tích ở những bước mà kỹ thuật
đó can thiệp~\autocite{p22whichlime}. Cột chữ cái là chỗ đọc được nhiều nhất, vì
nó cho biết cả lĩnh vực đã dồn công vào đâu.

Đếm theo nhóm vấn đề, với quy ước một kỹ thuật mang hai chữ cái thì được tính
vào cả hai nhóm: 23 kỹ thuật nhắm vào độ khớp, 20 vào tính ổn định, 13 vào khả
năng diễn giải, 11 vào tính cục bộ, và 3 vào hiệu năng. Tổng các con số ấy là 70
lượt gán trên 48 kỹ thuật, vì 20 kỹ thuật mang hai chữ cái và một kỹ thuật mang
ba; 27 kỹ thuật còn lại mang đúng một~\autocite{p22whichlime}.

Hai con số đầu đáng dừng lại. Độ khớp và tính ổn định gộp lại chiếm 43 trong 70
lượt gán, và nếu đếm theo kỹ thuật thay vì theo lượt gán thì đúng 32 trong 48 kỹ
thuật, tức hai phần ba, nhắm vào ít nhất một trong hai nhóm ấy. Hai nhóm ấy là
\tn{độ trung thực}{faithfulness} của mô hình thay thế và điều kiện cần của nó.
Nói cách khác, hai phần ba công sức của cả nhánh nghiên cứu đổ vào đúng chỗ mà
chương~\ref{ch:lime} đã chỉ ra là chưa được xác nhận, nên có thể nói lĩnh vực
biết LIME có vấn đề ở đó và biết từ lâu.

Con số 3 kỹ thuật nhắm vào hiệu năng thì nói điều ngược lại. Chi phí tính toán
là nhóm vấn đề duy nhất trong năm nhóm mà ai cũng đo được bằng một chiếc đồng
hồ, không cần tới dụng cụ đo nào đang tranh cãi, thế mà nó lại là nhóm ít người
làm nhất. Bốn nhóm đông người làm hơn đều là những nhóm mà việc xác nhận đã sửa
được hay chưa còn phải chờ một thước đo.

\index{độ trung thực!của mô hình thay thế}%
Con số 21 kỹ thuật mang từ hai chữ cái trở lên xác nhận điều nói ở cuối
mục~\ref{sec:ch10-nam-nhom-van-de}, lần này từ phía dữ liệu. Cặp hay gặp nhất là tính ổn
định đi với độ khớp, 10 kỹ thuật, và điều đó hợp lý theo đúng cơ chế: hai lần
chạy cho hai mô hình thay thế khác nhau thì hai mô hình ấy không thể cùng là
lời giải thích đúng cho cùng một quyết định, nên chữa được tính ổn định là điều
kiện cần để nói gì về độ khớp. Kỹ thuật
duy nhất mang ba chữ cái là B-LIME, nhắm cùng lúc vào tính ổn định, độ khớp và
khả năng diễn giải~\autocite{p22whichlime}.

Còn một trục nữa mà bảng ghi lại và mục này không đếm: các dấu tích ở bốn cột
bước. Chúng đọc được bằng mắt trên trang in nhưng không rút ra được đủ tin cậy
để in một con số, nên chương này không in con số nào từ chúng. Điều rút được mà
không cần đếm chính xác thì đã nằm ở mục~\ref{sec:ch10-bon-buoc}: có kỹ thuật
sửa một bước, có kỹ thuật sửa cả bốn.

Khảo sát còn một bảng thứ hai, xếp cùng 48 kỹ thuật ấy theo kiểu dữ liệu và
lĩnh vực ứng dụng: văn bản, bảng, ảnh, chuỗi thời gian, âm thanh, đồ thị, và
cặp ảnh kèm văn bản; bên trong mỗi kiểu lại tách ra những kỹ thuật chỉ được thử
trong một lĩnh vực hẹp như khảo cổ hay y tế~\autocite{p22whichlime}. Bảng ấy có
một dòng cuối đáng chép lại nguyên ý, vì nó là chỗ khảo sát bắt đầu quay lại
chất vấn chính cái nó vừa dựng: các tác giả nói họ không xác định dứt khoát
được rằng lời khẳng định về tính phổ dụng của từng kỹ thuật có đúng hay không,
vì không phải kỹ thuật nào cũng có mã nguồn~\autocite{p22whichlime}. Mục sau
theo đúng chỗ hở đó.
